%% kapa: exam=1 round=21 year=2010 period=2 form=A subject=appraisal_law subject_ko=부동산관계법규 seq=15 number=55
%% source: Q-Net 감정평가사 기출문제 (cst00309) artlSeq=5024181 file=extracted/1008/5024181/2010년 제21회 감정평가사 1차시험 기출문제/2교시.hwp
%% conversion: kordoc 4.24.51 -> tools/gen1.py (auto); review=proofread-gpt-sol; flags=none
\begin{question}{55}
국토의 계획 및 이용에 관한 법령상 장기미집행 도시계획시설부지의 매수 청구에 관한 다음 규정 중 ( ) 안에 들어갈 옳은 것을 순서대로 나열한 것은?
\begin{bogi}[다 음]
도시계획시설에 대한 도시관리계획의 결정(이하 "도시계획시설결정"이라 한다)의 고시일부터 ( ㄱ ) 이내에 그 도시계획시설의 설치에 관한 도시계획시설사업이 시행되지 아니하는 경우(제88조에 따른 실시계획의 인가나 그에 상당하는 절차가 진행된 경우는 제외한다. 이하 같다) 그 도시계획시설의 부지로 되어 있는 토지 중 지목(地目)이 ( ㄴ )인 토지(그 토지에 있는 건축물 및 정착물을 포함한다. 이하 이 조에서 같다)의 소유자는 대통령령으로 정하는 바에 따라 특별시장ㆍ광역시장ㆍ시장 또는 군수에게 그 토지의 매수를 청구할 수 있다. \par
\end{bogi}
\par\noindent\hspace{2em}\uline{\makebox[10em][l]{ㄱ}}\quad\uline{\makebox[10em][l]{ㄴ}}\par
\begin{choices}
  \item \makebox[10em][l]{20년}\quad 대(垈)
  \item \makebox[10em][l]{10년}\quad 대(垈)
  \item \makebox[10em][l]{20년}\quad 잡종지
  \item \makebox[10em][l]{10년}\quad 잡종지
  \item \makebox[10em][l]{20년}\quad 답(畓)
\end{choices}
\end{question}
\answer{2}
